\Needspace{8\baselineskip}\originalchapter{轴对称平衡与安全因子}\label{ch:gs}
\originalsection{磁通函数的几何意义}
柱坐标 $(R,\varphi,Z)$ 轴对称, 用每弧度极向磁通 $\Psi$ 定义
\begin{equation}
 \mathbf B=\frac1R\nabla\Psi\times\mathbf e_\varphi+\frac{F}{R}\mathbf e_\varphi,
 \quad B_R=-\frac{\Psi_Z}R,\quad B_Z=\frac{\Psi_R}R.
\end{equation}
这样自动满足无散条件且 $\mathbf B\cdot\nabla\Psi=0$. 平衡沿场分量给 $\mathbf B\cdot\nabla p=0$, 所以嵌套磁面中 p=p(\ensuremath{\Psi}). 环向力平衡又给 F=F(\ensuremath{\Psi}). 这两个通量函数要由热力、加热或电流约束指定, 不由几何单独决定.

安培方程的环向分量为 $\mu_0J_\varphi=-\Delta^*\Psi/R$, $\Delta^*=R\partial_R(R^{-1}\partial_R)+\partial_Z^2$. 将极向 J 与 B 代入剩余力平衡得
\begin{equation}\label{eq:gs}
 \Delta^*\Psi=-\mu_0R^2\frac{\dd p}{\dd\Psi}-F\frac{\dd F}{\dd\Psi}.
\end{equation}
这是 Grad--Shafranov 方程. 若通量采用总磁通 $2\pi\Psi$, 所有导数和系数必须整体换算, 不可只给压力项加2. 本册以此约定检查作业和考试公式.

\originalsection{磁力线绕行与 q}
磁力线绕极向一周积累的环向角除 $2\pi$ 定义安全因子
\begin{equation}
 q(\Psi)=\frac{F}{2\pi}\oint_\Psi\frac{\dd\ell_p}{R|\nabla\Psi|}.
\end{equation}
大长径比圆截面近似 $q=rB_\varphi/(R_0B_\theta)$. 若 $B_\theta=\mu_0I(r)/(2\pi r)$, 可从实测电流剖面算 q. q 本身不是``安全保证'', 而是绕行比; $q=m/n$ 时扰动 $\exp[\ii(m\theta-n\varphi)]$ 沿磁力线相位近似不变, 形成共振面. 磁剪切 $s=(r/q)q'$ 描述邻近面的绕行差, 对局部稳定性至关重要.

\originalsection{欧姆电流与模型闭合}
环向驱动 E、强平行传导下近似 $J_\parallel=E_\parallel/\eta(T_e)$. 对给定 T(r) 可用安培方程求 B\ensuremath{\theta} 再算 q, 但 T 又受欧姆热 $\eta J^2$ 与输运决定. 只给外部总电流还不足以确定 q(r). 理想平衡不是加热输运方程, 二者要以一致源项联立或明确只讨论指定剖面.
